\begin{figure}[tb]
    \newcommand{\subfigwidth}{0.49\linewidth} % Define a variable for the width
    \centering

    \begin{subfigure}[t]{\subfigwidth}
        \centering
        \includegraphics[width=\linewidth]{chapters_WACV/figures/scatterplots/light-blue/scatter_plot_DSO-1_VIPP_OpenForensics_IMD2020_random_distractions_QF_(85-95)_small_large_TruFor__WACV_square.pdf}
        \caption{Unmodified TruFor}
        \label{figure:motivation:TruFor-scatter}
    \end{subfigure}
    \hfill
    \begin{subfigure}[t]{\subfigwidth}
        \centering
        \includegraphics[width=\linewidth]{chapters_WACV/figures/scatterplots/light-blue/scatter_plot_DSO-1_VIPP_OpenForensics_IMD2020_random_distractions_QF_(85-95)_small_large_TruFor_maskedcnnManual_Automatic__WACV-dilated-5_square.pdf}
        \caption{Masked Convolution (ours)}
        \label{figure:evaluation:TruFor-scatter}
    \end{subfigure}

    \caption{
    Scatter plots of the F1* scores~\eqref{eq:F1*-score} of unmodified TruFor~\cite{TruFor} (a) vs. TruFor with the proposed masked-convolution approach (b), both with manual and automatic masking.
    Points farther to the left of the diagonal indicate stronger performance degradation due to forensic distractions.
    While the unmodified TruFor shows many points to the left, the adapted TruFor (both manual and automatic) shows a clear shift toward the diagonal, reflecting improved resilience to distractions. Notably, the points for manual masking (orange) cluster closely around the diagonal, indicating very high resilience to distractions.}
    \label{figure:motivation+evaluation:TruFor-scatter}
\end{figure}

